\documentclass[11pt,a4paper]{article}

\usepackage[margin=1in]{geometry}
\usepackage[T1]{fontenc}
\usepackage[utf8]{inputenc}
\usepackage{lmodern}
\usepackage{microtype}
\usepackage{amsmath}
\usepackage{amssymb}
\usepackage{booktabs}
\usepackage{array}
\usepackage{longtable}
\usepackage{hyperref}
\usepackage{xcolor}

\hypersetup{
  colorlinks=true,
  linkcolor=blue!50!black,
  urlcolor=blue!50!black,
  pdftitle={Raman Transition Calculator Manual},
  pdfauthor={OpenAI Codex},
  pdfsubject={Python translation of Raman.txt},
  pdfkeywords={Raman transition, atomic interferometer, Rabi oscillation, Python, Tkinter}
}

\setlength{\parskip}{0.6em}
\setlength{\parindent}{0pt}

\title{Raman Transition Calculator Manual\\\large Python Translation of \texttt{Raman.txt}}
\author{Project workspace: \texttt{Raman\_calculator}}
\date{\today}

\begin{document}

\maketitle
\tableofcontents

\section{Purpose}

This document describes the Python implementation of the original Mathematica script
\texttt{Raman.txt}. The translated application computes Raman-driven Rabi oscillations in an
atomic interferometer, displays the free-expansion growth of the transverse atom-cloud size, and
integrates a second Raman-detuning workflow translated from the user's standalone Python tool in
\texttt{Raman deturning/Raman\_deturning\_4.0.py}.

The implementation goals are:

\begin{itemize}
  \item preserve the structure and default constants of the Mathematica model,
  \item preserve the calculation logic of the original Raman-detuning script,
  \item expose the main experimental parameters through a clean English-language graphical interface,
  \item provide direct plots for both transition probability and cloud-size evolution,
  \item provide a dedicated calibration workflow for reconstructing $\alpha$ and $v_x$ from
  flying-up and falling-down detuning scans,
  \item document the assumptions, units, and numerical method in a form suitable for research use.
\end{itemize}

\section{Source-to-Python Translation}

The Mathematica source defines a model with the following components:

\begin{itemize}
  \item Gaussian Raman-beam intensities,
  \item an effective Raman coupling frequency depending on the large single-photon detuning,
  \item a Gaussian longitudinal velocity distribution,
  \item a Gaussian transverse spatial distribution that broadens during free expansion,
  \item a two-dimensional average of the Raman transition probability over radius and velocity.
\end{itemize}

The Python code preserves these elements in \texttt{raman\_model.py}. The graphical interface in
\texttt{app.py} provides controlled access to the same quantities and also reports derived values
such as the on-axis Raman Rabi frequency and the estimated $\pi$-pulse time.

The project now also contains an integrated detuning model in \texttt{raman\_detuning\_model.py}.
That module preserves the branch logic of the earlier standalone Python detuning calculator and
adds a calibration routine that recovers the Raman-beam angle $\alpha$ and transverse velocity
$v_x$ from two signed detuning measurements.

\section{Default Parameters}

The table below lists the translated defaults inherited from the Mathematica file.

\begin{center}
\begin{tabular}{>{\raggedright\arraybackslash}p{0.52\linewidth}r}
\toprule
Quantity & Default value \\
\midrule
Recoil velocity $v_{\mathrm{rec}}$ & $5.88 \times 10^{-3}$ m/s \\
Initial cloud size $\sigma_r(0)$ & 5 mm \\
Large Raman detuning $\Delta / 2\pi$ & $-1000$ MHz \\
Fine-structure offset $\delta_2 / 2\pi$ & 157 MHz \\
Natural linewidth $\gamma / 2\pi$ & 6.065 MHz \\
Raman beam power $P_1$ & 14.0 mW \\
Raman beam power $P_2$ & 5.384615 mW \\
Beam waist $w_0$ & 11.5 mm \\
Free-expansion time $T$ in $P(T,\tau)$ & 56 ms \\
Attenuation factor & 1 \\
Coupling gain $G$ & 1 \\
\bottomrule
\end{tabular}
\end{center}

\subsection{Temperature Mapping}

The original Mathematica file directly specifies the transverse velocity spread as
$\sigma_{v,T} = 2.65\,v_{\mathrm{rec}}$. The Python application uses the transverse temperature as
the primary thermal input for that branch of the model. Therefore, the default transverse
temperature is chosen such that

\begin{equation}
\sqrt{\frac{k_B T}{M}} = 2.65\,v_{\mathrm{rec}},
\end{equation}

which gives

\begin{equation}
T_{\mathrm{default}} = 2.537918~\mu\mathrm{K}.
\end{equation}

The GUI distinguishes two thermal directions:

\begin{itemize}
  \item the \textbf{transverse temperature}, which sets $\sigma_{v,T}$ and therefore the radial
  beam convolution and cloud expansion,
  \item the \textbf{longitudinal temperature}, which sets $\sigma_{v,L}$ and therefore the
  velocity-selection distribution $f(v_L)$.
\end{itemize}

By default, the longitudinal temperature is linked to the transverse temperature so that the two
directions start from the same value. The longitudinal value can then be released and edited
independently whenever the experiment requires different thermal widths along and across the Raman
beam axis.

\section{Physical Model}

\subsection{Gaussian Beam Intensities}

The Gaussian intensity of each Raman beam is

\begin{equation}
I_i(r) = \frac{2P_i}{\pi w_0^2}\exp\!\left(-\frac{2r^2}{w_0^2}\right),
\end{equation}

where $i \in \{1,2\}$, $P_i$ is the optical power, and $w_0$ is the common beam waist.

\subsection{Effective Raman Coupling}

The translated effective Raman Rabi frequency is

\begin{equation}
\Omega_{\mathrm{eff}}(r)
=
\frac{\gamma^2 G}{\mathrm{attn}}
\frac{\sqrt{I_1(r)I_2(r)}}{2I_{\mathrm{sat}}}
\left[
\frac{5}{24\Delta} + \frac{3}{24(\Delta+\delta_2)}
\right]\frac{1}{2},
\end{equation}

where:

\begin{itemize}
  \item $\Delta$ is the large Raman detuning (\texttt{desacc} in the code),
  \item $\delta_2$ is the offset corresponding to the additional excited-state contribution,
  \item $I_{\mathrm{sat}} = 16.7$ W/m$^2$,
  \item $\mathrm{attn}$ is the attenuation factor from the original script,
  \item $G$ is a dimensionless gain term retained from the Mathematica implementation.
\end{itemize}

In the Python code the absolute value of the resulting coupling is used when reporting an effective
frequency, because the transition probability depends on $\Omega_{\mathrm{eff}}^2$.

\subsection{Velocity and Spatial Distributions}

The longitudinal velocity distribution is Gaussian:

\begin{equation}
f(v_L) = \frac{1}{\sqrt{2\pi}\,\sigma_{v,L}}
\exp\!\left(-\frac{v_L^2}{2\sigma_{v,L}^2}\right).
\end{equation}

The transverse cloud size after free-expansion time $T$ is

\begin{equation}
\sigma_r(T) = \sqrt{\sigma_r(0)^2 + \sigma_{v,T}^2 T^2}.
\end{equation}

The associated radial probability density is

\begin{equation}
h(r,T) = \frac{1}{2\pi \sigma_r(T)^2}
\exp\!\left(-\frac{r^2}{2\sigma_r(T)^2}\right).
\end{equation}

\subsection{Generalized Rabi Frequency}

For each $(r, v_L)$ point, the generalized Raman Rabi frequency is

\begin{equation}
\Omega_R(r,v_L) =
\sqrt{\Omega_{\mathrm{eff}}(r)^2 + \left(k_{\mathrm{eff}}v_L - d\right)^2},
\end{equation}

with

\begin{equation}
k_{\mathrm{eff}} = \frac{4\pi}{\lambda},
\end{equation}

and $d$ the user-specified two-photon detuning.

\subsection{Transition Probability}

The Raman transition probability is evaluated as

\begin{equation}
P(T,\tau)
=
\int \int
2\pi r \, f(v_L)\, h(r,T)\,
\left[
\frac{\Omega_{\mathrm{eff}}(r)}{\Omega_R(r,v_L)}
\sin\!\left(\frac{\Omega_R(r,v_L)\tau}{2}\right)
\right]^2
\mathrm{d}v_L\,\mathrm{d}r.
\end{equation}

This is a direct Python translation of the Mathematica function
\texttt{Ptrans[T\_, tau\_, d\_, w0\_, attn\_]}.

\section{Raman Detuning Model}

\subsection{Wave-vector and Recoil Definitions}

The integrated detuning tool uses the same conventions as the original standalone Python script.
The laser wave number is

\begin{equation}
k = \frac{2\pi}{\lambda},
\end{equation}

the effective Raman wave number is

\begin{equation}
k_{\mathrm{eff}} = 2k = \frac{4\pi}{\lambda},
\end{equation}

and the recoil term is represented as an angular frequency

\begin{equation}
\omega_r = 2\pi f_r,
\end{equation}

where $f_r$ is entered in kHz in the GUI.

\subsection{Four Detuning Branches}

Let $v_z$ be the known vertical atomic velocity, $\alpha$ the Raman-beam angle relative to the
$z$ axis, and $v_x$ the transverse velocity to be inferred or propagated. The transition direction
controls the recoil sign through

\begin{equation}
s_{\mathrm{tr}} =
\begin{cases}
+1, & \text{for } \mathrm{F1}\rightarrow\mathrm{F2}, \\
-1, & \text{for } \mathrm{F2}\rightarrow\mathrm{F1}.
\end{cases}
\end{equation}

The detuning model then provides four signed branches:

\begin{align}
\delta_{\mathrm{up},+} &= k_{\mathrm{eff}}v_z\sin\alpha + k_{\mathrm{eff}}v_x\cos\alpha + s_{\mathrm{tr}}\omega_r, \\
\delta_{\mathrm{up},-} &= -k_{\mathrm{eff}}v_z\sin\alpha - k_{\mathrm{eff}}v_x\cos\alpha + s_{\mathrm{tr}}\omega_r, \\
\delta_{\mathrm{down},+} &= k_{\mathrm{eff}}v_z\sin\alpha - k_{\mathrm{eff}}v_x\cos\alpha + s_{\mathrm{tr}}\omega_r, \\
\delta_{\mathrm{down},-} &= -k_{\mathrm{eff}}v_z\sin\alpha + k_{\mathrm{eff}}v_x\cos\alpha + s_{\mathrm{tr}}\omega_r.
\end{align}

In the GUI these four cases are reported as:

\begin{itemize}
  \item \texttt{flying up, $\Delta>0$},
  \item \texttt{flying up, $\Delta<0$},
  \item \texttt{falling down, $\Delta>0$},
  \item \texttt{falling down, $\Delta<0$}.
\end{itemize}

The sign of the entered or computed detuning therefore matters operationally: it is used to
identify which branch of the original model is active.

\subsection{Velocity Inversion}

When the user selects the \texttt{Detuning -> vx} mode, the application keeps the original branch
logic of the standalone script. The user provides:

\begin{itemize}
  \item a signed detuning in kHz,
  \item the motion direction (\texttt{flying up} or \texttt{falling down}),
  \item the transition direction (\texttt{F1→F2} or \texttt{F2→F1}).
\end{itemize}

The sign of the detuning automatically selects either the $\Delta>0$ or $\Delta<0$ branch, and
the program returns both the recovered $v_x$ and the branch that was used internally.

\subsection{Calibration from Flying-up and Falling-down Scans}

The calibration page addresses the case in which $\alpha$ and $v_x$ are both initially unknown.
The user enters:

\begin{itemize}
  \item the signed detuning extracted from the \texttt{flying up} scan,
  \item the signed detuning extracted from the \texttt{falling down} scan,
  \item the transition direction used in both scans.
\end{itemize}

Define the sign selectors

\begin{equation}
s_{\mathrm{up}} = \mathrm{sign}(\delta_{\mathrm{up}}), \qquad
s_{\mathrm{down}} = \mathrm{sign}(\delta_{\mathrm{down}}),
\end{equation}

and the branch-reduced quantities

\begin{equation}
U = \frac{\delta_{\mathrm{up}} - s_{\mathrm{tr}}\omega_r}{s_{\mathrm{up}}}, \qquad
D = \frac{\delta_{\mathrm{down}} - s_{\mathrm{tr}}\omega_r}{s_{\mathrm{down}}}.
\end{equation}

These satisfy

\begin{equation}
U = k_{\mathrm{eff}}v_z\sin\alpha + k_{\mathrm{eff}}v_x\cos\alpha,
\end{equation}

\begin{equation}
D = k_{\mathrm{eff}}v_z\sin\alpha - k_{\mathrm{eff}}v_x\cos\alpha.
\end{equation}

Therefore, the calibration step reconstructs the beam angle and transverse velocity from

\begin{equation}
\sin\alpha = \frac{U + D}{2k_{\mathrm{eff}}v_z},
\end{equation}

\begin{equation}
v_x = \frac{U - D}{2k_{\mathrm{eff}}\cos\alpha}.
\end{equation}

This is exactly the logic implemented in the integrated Python calibration routine. After a
successful calibration, the GUI can apply the recovered $\alpha$ back into the detuning constants
and copy the recovered $v_x$ into the calculator input field for immediate reuse.

\section{Graphical Interface}

The desktop interface is implemented with Tkinter and Matplotlib. The GUI is organized into two
top-level pages, \texttt{Simulation} and \texttt{Detuning}, so that the full Raman-transition and
detuning workflows remain available in a single application.

\subsection{Page Structure}

\begin{itemize}
  \item \texttt{Simulation}: Raman transition probability, atom-cloud expansion, presets, locked
  overlays, and interactive plotting.
  \item \texttt{Detuning}: Raman detuning forward calculation, inverse $v_x$ recovery, and
  calibration of $\alpha$ and $v_x$ from scan results.
\end{itemize}

\subsection{Preset Library}

The interface now includes a preset library for rapid switching between common Raman operating
conditions. Four presets are bundled with the application:

\begin{center}
\begin{tabular}{>{\raggedright\arraybackslash}p{0.25\linewidth} >{\raggedright\arraybackslash}p{0.62\linewidth}}
\toprule
Preset & Applied values \\
\midrule
\texttt{Original Raman.txt Defaults} & Direct translation of the Mathematica defaults. \\
\texttt{Raman Down} & $P_1 = 14$ mW, $P_2 = 7$ mW, $w_0 = 11.5$ mm, $T = 56$ ms, $\Delta / 2\pi = -1300$ MHz. \\
\texttt{Raman Up} & $P_1 = 150$ mW, $P_2 = 70$ mW, $w_0 = 30$ mm, $T = 78$ ms, all remaining parameters inherited from the default set. \\
\texttt{Raman Labeling} & $P_1 = 150$ mW, $P_2 = 70$ mW, $w_0 = 30$ mm, $T = 780$ ms, all remaining parameters inherited from the default set. \\
\bottomrule
\end{tabular}
\end{center}

The preset values labelled by the user as ``tau = 56 ms'', ``tau = 78 ms'', and ``tau = 780 ms''
are implemented as the advanced expansion-time parameter $T$, because the plotted Raman pulse
duration $\tau$ remains the horizontal-axis sweep in $\mu$s.

The preset panel supports four actions:

\begin{itemize}
  \item apply the currently selected preset and recompute the plots,
  \item save the current parameter set into the selected preset name,
  \item save the current parameter set as a new preset,
  \item delete a locally stored preset or a local override of a bundled preset.
\end{itemize}

User-created presets and local overrides are stored in \texttt{presets.json}.

\subsection{Simulation-page Action Bar}

The interface now places the main execution controls in the top-right corner of the application
window whenever the \texttt{Simulation} page is active. This keeps the numerical input column
focused on parameters while reserving a dedicated location for simulation management:

\begin{itemize}
  \item \texttt{Run Simulation}: execute the current parameter set,
  \item \texttt{Reset to Defaults}: restore the translated \texttt{Raman.txt} defaults,
  \item \texttt{Export Current CSV}: export the currently active simulation only,
  \item \texttt{Lock Current}: preserve the current curves as reference traces,
  \item \texttt{Clear Locked}: remove all locked reference traces.
\end{itemize}

\subsection{Primary Inputs}

\begin{longtable}{>{\raggedright\arraybackslash}p{0.30\linewidth} >{\raggedright\arraybackslash}p{0.15\linewidth} >{\raggedright\arraybackslash}p{0.45\linewidth}}
\toprule
Field & Unit & Meaning \\
\midrule
Transverse temperature & $\mu$K & Thermal input used to define the transverse velocity spread $\sigma_{v,T}$, the radial beam convolution, and the cloud expansion. \\
Use a separate longitudinal temperature & boolean & Unlocks an independent longitudinal thermal input. When disabled, the longitudinal temperature is linked to the transverse value. \\
Longitudinal temperature & $\mu$K & Thermal input used to define the longitudinal velocity spread $\sigma_{v,L}$ entering the velocity-selection distribution $f(v_L)$. \\
Large Raman detuning desacc / $2\pi$ & MHz & One-photon detuning entering the effective Raman coupling. \\
Beam power $P_1$ & mW & Optical power of Raman beam 1. \\
Beam power $P_2$ & mW & Optical power of Raman beam 2. \\
Beam waist $w_0$ & mm & Common Gaussian waist used in the intensity model. \\
$\tau$ minimum & $\mu$s & Lower bound of the pulse-duration sweep. \\
$\tau$ maximum & $\mu$s & Upper bound of the pulse-duration sweep. \\
Number of $\tau$ samples & points & Plot resolution for the pulse-duration sweep. \\
\bottomrule
\end{longtable}

\subsection{Advanced Model Parameters}

\begin{longtable}{>{\raggedright\arraybackslash}p{0.30\linewidth} >{\raggedright\arraybackslash}p{0.15\linewidth} >{\raggedright\arraybackslash}p{0.45\linewidth}}
\toprule
Field & Unit & Meaning \\
\midrule
Expansion time $T$ before Raman pulse & ms & Free-expansion time used inside the Raman-transition integral. \\
Initial cloud size $\sigma_r(0)$ & mm & Initial transverse root-mean-square cloud size. \\
Two-photon detuning $d / 2\pi$ & kHz & Residual Raman detuning used in the generalized Rabi frequency. \\
Attenuation factor & a.u. & Multiplicative suppression of $\Omega_{\mathrm{eff}}$. \\
Coupling gain $G$ & a.u. & Dimensionless gain factor retained from the original source. \\
\bottomrule
\end{longtable}

\subsection{Numerical Controls}

\begin{longtable}{>{\raggedright\arraybackslash}p{0.30\linewidth} >{\raggedright\arraybackslash}p{0.15\linewidth} >{\raggedright\arraybackslash}p{0.45\linewidth}}
\toprule
Field & Unit & Meaning \\
\midrule
Radial grid points & points & Number of quadrature points in the radial integral. \\
Velocity grid points & points & Number of quadrature points in the longitudinal velocity integral. \\
Radial cutoff & $w_0$ units & Upper radial integration limit in multiples of the beam waist. \\
Velocity cutoff & $\sigma_{v,L}$ units & Longitudinal velocity truncation used in the numerical integral. \\
\bottomrule
\end{longtable}

\subsection{Detuning-page Inputs}

\begin{longtable}{>{\raggedright\arraybackslash}p{0.30\linewidth} >{\raggedright\arraybackslash}p{0.15\linewidth} >{\raggedright\arraybackslash}p{0.45\linewidth}}
\toprule
Field & Unit & Meaning \\
\midrule
\texttt{vz} & m/s & Known vertical atomic velocity used in the detuning model. \\
\texttt{alpha} & deg & Raman-beam angle relative to the vertical axis. Used by the forward and inverse detuning calculators; not required as an input for calibration. \\
Laser wavelength & nm & Optical wavelength used to construct $k$ and $k_{\mathrm{eff}}$. \\
Recoil frequency & kHz & Two-photon recoil-frequency term converted internally to $\omega_r$. \\
Mode selector & --- & Switches between \texttt{vx -> Detuning} and \texttt{Detuning -> vx}. \\
Input field & mm/s or kHz & Accepts either $v_x$ or a signed detuning, depending on the selected mode. \\
Motion direction & --- & Specifies \texttt{flying up} or \texttt{falling down} for inverse calculations. \\
Transition direction & --- & Specifies \texttt{F1→F2} or \texttt{F2→F1}. \\
Calibration flying-up detuning & kHz & Signed result extracted from the flying-up scan. \\
Calibration falling-down detuning & kHz & Signed result extracted from the falling-down scan. \\
Calibration transition & --- & Transition direction used during the calibration scans. \\
\bottomrule
\end{longtable}

\section{Numerical Method}

The original Mathematica script used adaptive Monte Carlo integration. In the Python translation, a
deterministic tensor-product quadrature is used instead:

\begin{itemize}
  \item the radial coordinate is sampled on a uniform grid from $0$ to the user-selected radial cutoff,
  \item the longitudinal velocity is sampled on a uniform grid from
  $-\alpha \sigma_{v,L}$ to $+\alpha \sigma_{v,L}$, where $\alpha$ is the chosen cutoff multiplier,
  \item the integrals are evaluated using trapezoidal quadrature,
  \item the pulse-duration array is processed in chunks to keep memory usage modest while remaining fast enough for interactive use.
\end{itemize}

This method is appropriate for a responsive research tool, because the integrand is smooth and the
user can directly increase the grid resolution whenever additional convergence margin is required.

\section{Outputs and Interpretation}

The application reports two primary plots.

\subsection{Rabi Oscillation Plot}

The first plot displays the transition probability $P(T,\tau)$ over the selected pulse-duration
range. A dashed vertical line marks the estimated on-axis $\pi$-pulse duration whenever it lies
inside the plotted window. The plot tab includes an automatic vertical-axis scaling option so that
small-amplitude Raman signals are displayed with an appropriate zoom level rather than being forced
onto a fixed $0$ to $1$ scale.

\subsection{Cloud-Size Evolution Plot}

The second plot displays

\begin{equation}
\sigma_r(t) = \sqrt{\sigma_r(0)^2 + \sigma_{v,T}^2 t^2}
\end{equation}

from $t = 0$ to the user-selected free-expansion time $T$. The horizontal axis is displayed in an
automatically chosen unit (\textmu s, ms, or s) depending on the selected expansion interval.

This behavior is intentionally aligned with the original Mathematica definition
\texttt{dispPos[T]}, which describes cloud broadening during free expansion prior to the Raman
interaction. The same fixed interaction time $T$ is then used inside the Raman-transition integral
$P(T,\tau)$.

\subsection{Derived Values Panel}

The derived-values panel summarizes:

\begin{itemize}
  \item longitudinal and transverse velocity spreads,
  \item initial and expanded cloud sizes,
  \item on-axis peak beam intensities,
  \item on-axis effective Raman Rabi frequency,
  \item estimated on-axis $\pi$-pulse time,
  \item numerical grid settings and integration cutoffs.
\end{itemize}

\subsection{Interactive Plot Tools}

The Matplotlib toolbar remains available for pan and zoom. In addition, the application provides
direct in-figure interaction:

\begin{itemize}
  \item hover readout of the nearest sampled point on either curve,
  \item click-to-pin markers with numerical coordinate labels,
  \item double-click or the \texttt{Reset Plot View} button to restore the latest auto-generated axis limits,
  \item a \texttt{Clear Plot Markers} button to remove pinned annotations.
\end{itemize}

\subsection{Locked Result Overlays}

The application can preserve the current simulation result as a locked reference while new
parameter sets are computed. When multiple simulations are displayed simultaneously:

\begin{itemize}
  \item locked curves are drawn as dashed traces,
  \item the current simulation is drawn as a solid trace,
  \item plot legends are enabled automatically,
  \item each legend label is generated from only the parameters that differ across the displayed
  simulations, keeping the legend concise while still informative.
\end{itemize}

\subsection{Detuning-page Outputs}

The \texttt{Detuning} page provides three text-based result areas:

\begin{itemize}
  \item \texttt{Calculator Result}: reports either the four detuning branches or the recovered
  $v_x$ and the internally selected branch,
  \item \texttt{Calibration Result}: reports the reconstructed $\alpha$, reconstructed $v_x$, and
  the sign branch identified for both the flying-up and falling-down scan inputs,
  \item \texttt{Detuning Notes}: summarizes the model scope and the meaning of the sign
  conventions.
\end{itemize}

\section{How to Run the Software}

From the project directory:

\begin{verbatim}
python3 app.py
\end{verbatim}

If dependencies need to be installed:

\begin{verbatim}
python3 -m pip install -r requirements.txt
\end{verbatim}

To execute the smoke tests:

\begin{verbatim}
python3 -m unittest discover -s tests -v
\end{verbatim}

To rebuild this PDF:

\begin{verbatim}
pdflatex -interaction=nonstopmode -halt-on-error \
  -output-directory=docs docs/raman_calculator_manual.tex
\end{verbatim}

\section{Project File Summary}

\begin{center}
\begin{tabular}{>{\raggedright\arraybackslash}p{0.36\linewidth} >{\raggedright\arraybackslash}p{0.52\linewidth}}
\toprule
File & Role \\
\midrule
\texttt{Raman.txt} & Original Mathematica source retained for reference. \\
\texttt{raman\_model.py} & Physics model, unit conversions, validation, and numerical solver. \\
\texttt{app.py} & Graphical desktop application with plots and parameter entry. \\
\texttt{presets.json} & User-written preset store containing custom presets and local overrides. \\
\texttt{tests/test\_raman\_model.py} & Numerical smoke tests for the translated model. \\
\texttt{README.md} & Quick-start and project summary. \\
\texttt{docs/raman\_calculator\_manual.tex} & Full technical manual. \\
\bottomrule
\end{tabular}
\end{center}

\section{Closing Remarks}

This Python translation is intentionally conservative: it stays close to the original Mathematica
structure, makes units explicit, and separates the physical model from the interface so that future
extensions remain straightforward. Typical next steps for a laboratory workflow would be to add
batch export of parameter scans, fit measured Rabi data, or expose independent longitudinal and
transverse temperatures if the experiment requires a more anisotropic thermal model.

\end{document}
